\SourceLeaf{530}which contained a correction to Article 139 of the Disqu. Arithm., has been inserted there at the relevant passage. The note on page 91 is taken from a handwritten note. Apart from this, the present edition differs from the earlier ones only in the correction of some printing errors. My additions are identified by square brackets [\,].
\par The \textit{Table of the Quadratic Character of Primes} is printed in the manner of Tabula II of the Disqu. Arithm., described in Article 99 and especially used (in Article 331) for factoring numbers. The handwriting of the manuscript entitled ‘\textit{Quadratorum numeris primis divisorum residua lateralia}’ most closely resembles that of the second part of the table for converting common fractions into decimal fractions. In place of the horizontal strokes indicating a quadratic residue, it contains small circles; those in adjacent horizontal or vertical rows are always connected by lines. While proofreading I noticed several errors. I then found 190 discrepancies in the characters in a single comparison with \textsc{Jacobi}’s \textit{Canon Arithmeticus}, and by determining them directly found agreement with those printed tables. The edition here has been corrected accordingly.
\par The first part of the \textit{Table for Converting Common Fractions into Decimal Fractions} is arranged here similarly to Tabula III of the Disqu. Arithm. For the primes and their powers $p^\pi$ between 3 and 463 it contains the mantissas $(1)$, $(2)\dotstwo(0)$ of the decimal fractions of $\frac{10\pdot r}{p^\pi}$, $\frac{10\pdot rr}{p^\pi}\dotstwo\frac{10}{p^\pi}$. Here $r$ denotes unity, so that $(1)=(2)=\dotstwo(0)$, when 10 is a primitive root of $p^\pi$; otherwise $r$ denotes the smallest of those primitive roots of $p^\pi$ for which, as base, the index of 10 takes its smallest value. To make the table easier to use, I have added the values other than 1 of $r$ on page 420. The manuscript, which also does not yet give the distinguishing digits of the different periods, most closely resembles the Analysis residuorum in appearance and seems to precede in date the piece placed here as the second part of the whole table.
This second part contains, for the primes and their power $p^\pi$ between 467 and 997, the mantissas of the decimal fractions of $\frac{100}{p^\pi}$. The manuscript gives the divisors in decreasing order and ends with the words: \textit{Explicitus October} 11. 1795. In print I have added, for the divisor 191 in period $(1)$, the 71st digit, and omitted, for the divisor 829, a number between the 151st and 152nd digits.
\par The first part of the \textit{Table of the Frequency of Primes}, discussed by \textsc{Gauss} himself in a letter (page 444), giving the number of primes in each of the first 1000 thousand groups, exists in a manuscript by \textsc{Gauss}. His manuscript remains, however, do not contain the counts mentioned in the letter of the hundred groups within the first million that contain a specified number of primes. The other part of the table, for the second and third million, comes from a manuscript by \textsc{Goldschmidt} alone. Herr \textsc{Meissel}, by counting and by his formula, has found the following corrections to pages 436 and 437:
